\section{Effect of the Calibration Reference}
\label{sec:tail}

\subsection{Threshold displacement}

Excluding scans with an official consensus nodule removes 14 of the \PCal{} negative calibration patients and leaves \UCal{}. Three of the six largest maxima of the selected detector belong to excluded patients, and each of them matches an official nodule in 3D (Fig.~\ref{fig:tail}a). The original threshold is the 62nd score, a logit of \PairCutoffP{}, whereas the new threshold is the 49th of 53 scores, a logit of \PairCutoffU{} that ranked 55th in the original pool, a drop of \PairDrop{} ranks. Removing 14 patients at random, with the same number per calibration subset, moves the threshold by a median of \PairRobustMedian{} ranks and by at most \PairControlHi{} ranks in 95\% of the draws. The displacement is therefore driven by which patients are removed, not by the reduced pool size.

At the lower threshold, the same scores recover \PairHitsU{} of the 42 omissions instead of \PairHitsP{} (Fig.~\ref{fig:tail}b), and five of the ten additional detections were already ranked above every non-target candidate in their scan. Alerts on the same \UTest{} test patients increase from \PairAlertsP{} to \PairAlertsU{} (Fig.~\ref{fig:tail}c), which is consistent with the nominal rate of 10\%.

\subsection{Robustness}

We repeated the replay for all \NHeadRows{} fitted configurations, which yield \NHeads{} distinct score vectors over hand-crafted, multi-view and 3D features. The threshold falls for \LowerHeads{} of them, by a median of \MedianDrop{} ranks, and for \PastBand{} the drop exceeds the 95\% range of random removal. Since these score vectors share patients and include correlated training variants, the count indicates the breadth of the effect rather than independent replications.

Fig.~\ref{fig:robust}a and Tab.~\ref{tab:reference} detail eight representative detectors, listing omissions found (of 42), targets found (of 276) and alerts among the 90 common negative test patients. The multi-view features are the embedding of an earlier 2.5D development model (Sec.~E of the supplement), and ``full ref.'' denotes heads trained on the full reference labels. The thresholds of the descriptor heads move by two ranks, within the random range, and both heads retain their two omission detections. In contrast, each of the six learned heads shown moves by six or seven ranks, outside the random range. A lower threshold produces additional detections only when the representation places targets between the two thresholds: the 3D first-read and full-reference heads gain 7 and 9 omissions, respectively, but the mark IPW head gains only 3. Changing the reader pair that defines the negatives yields between one and five detections for the selected detector, with at most one alert among the \RoleTest{} patients that are negative under every pair and under the official list (Fig.~\ref{fig:robust}b). The multi-view and 3D first-read heads detect 11 and 6 omissions, respectively, with every pair except the original C/D.

\begin{table}[t]
\centering\small
\setlength{\tabcolsep}{3pt}
\caption{Counts under the C/D threshold $\to$ after excluding official positives, with the threshold drop in ranks of the original pool.}
\label{tab:reference}
\begin{tabular}{lrrrr}
\toprule
Scorer / training & Omit. & All & Alert & Drop \\
\midrule
Descriptors / first & 2 $\to$ 2 & 39 $\to$ 55 & 7 $\to$ 12 & 2 \\
Descriptors / full ref. & 2 $\to$ 2 & 47 $\to$ 57 & 7 $\to$ 10 & 2 \\
Multi-view / first & 4 $\to$ 11 & 66 $\to$ 125 & 1 $\to$ 7 & 6 \\
Multi-view / full ref. & 3 $\to$ 11 & 69 $\to$ 142 & 1 $\to$ 9 & 7 \\
3D / first & 4 $\to$ 11 & 81 $\to$ 133 & 0 $\to$ 5 & 6 \\
3D / full ref. & 1 $\to$ 10 & 74 $\to$ 140 & 0 $\to$ 8 & 7 \\
3D / mark IPW & 3 $\to$ 6 & 70 $\to$ 91 & 2 $\to$ 6 & 6 \\
3D / paired tail & 1 $\to$ 11 & 69 $\to$ 138 & 0 $\to$ 8 & 7 \\
\bottomrule
\end{tabular}

\end{table}

\subsection{Effect on the method comparison}

Averaged over the six matched draws, excluding official positives raises the omission recall of the selected method by \PairChange{} points \PairChangeCI{}. Its difference to mark IPW moves from $\PairRelativeP$ to $+\PairRelativeU$ points \PairRelativeUCI{}, a shift of \ReferenceInteraction{} points (computed before rounding) \ReferenceInteractionCI{} (Fig.~\ref{fig:effects}b). The same shift occurs against half-budget mark stratification with discriminative augmentation (Mark disc.) ($-4.4$ to $+4.4$ points), whereas the gap to first-read training remains negative. The gains in all-target recall increase under the new pool (Fig.~\ref{fig:effects}c), and the alert differences stay within about two points (Fig.~\ref{fig:effects}d), with realized alert rates of 7.2\% for paired-tail sampling and 6.1\% for mark IPW. The primary conclusion of Sec.~\ref{sec:locked} thus hinges on how the negative pool is defined rather than on the scores themselves.

\begin{figure}[t]
\centering
\includegraphics[width=\columnwidth]{ct_evidence}
\caption{CT examples (28\,mm field of view; circle: matching radius): (a)~the candidate setting the C/D threshold, (b)~a maximum on an official-positive scan, (c)~an omission found only under the lower threshold, (d)~an omission below both. Margins: best target logit minus each threshold.}
\label{fig:ct}
\end{figure}

Finally, Fig.~\ref{fig:ct} relates these ranks to images: a consensus-positive maximum (b) leaves the negative pool with its score unchanged, and a test omission (c) then exceeds the lower threshold with the same image, location and score. Other omissions, such as (d), remain far below both thresholds, which illustrates why ranking and calibration need to be assessed separately.
